\subsection{A charged ball in the modified metric}
\label{sec:modified-ball-charge}

This comparison tests a particular degeneration of the modified Yamabe
representatives when the spectral deficit tends to zero and a uniform
upper ball-volume ratio is available. The negative Euler trace below
restricts possible initial ends of that Ricci flow. The regularized
flow used for fixed-type attainment has its own compactness argument
\cite{Baker2026BourAttainment}.

\begin{proposition}\label{prop:modified-ball-charge}
In the fixed subcritical band of Section~\ref{sec:modified-ricci-concentration}, let the modified metrics have vanishing spectral deficit on nonstandard types and
suppose additionally that
\[
 \operatorname{Vol}_{q_i} B_{q_i}(x,r)\le V r^4
 \qquad(x\in M_i,\ 0<r\le r_0)
\]
for fixed \(r_0,V>0\). There are \(c>0\), centers \(x_i\), and
radii \(a_i\to0\)
for which
\[
 \int_{B_{q_i}(x_i,a_i)} |Z_{q_i}|^2\,dV_{q_i}\ge c.
\]
On the same balls the scalar and both Weyl square integrals vanish
in the limit, and \(a_i\sqrt{\max|\Rm_{q_i}|}\to\infty\).
\end{proposition}

\begin{proof}
Use the smooth conformal approximations in
Theorem~\ref{thm:modified-shrinking-limits}. They preserve the
volume ratio upper bound with fixed enlarged constants on a smaller
interval, and distort distances by \(1+o(1)\). They have scalar
curvature at least \(-1\), a common inhomogeneous Sobolev bound,
and negative parts of both isotropic pair sums tending uniformly
to zero. We first prove that for a fixed sufficiently small
\(\epsilon>0\) and each fixed \(0<r<\min(1,r_0/2)\), at all late
indices some radius-\(r\) ball of the smoothed metric has full
curvature \(L^2\) norm greater than \(\epsilon\).

Otherwise fix one such radius and rescale by \(r^{-2}\) on the
contrary subsequence. The initial ball volume ratios through unit
scale and Sobolev constants have upper bounds independent of \(r\).
The rescaled zero-order Sobolev coefficient is at most its original
value. Apply Lee--Lee's local smoothing theorem with center arbitrary
and domain parameter larger than the diameter of each rescaled
closed manifold \cite[Theorem~1.1]{LeeLee2026Local}.
The theorem's constants are independent of that domain parameter.
It supplies a closed flow up to a common time \(t_*>0\), which we
take at most \(\min(1,\epsilon^2)\). In dimension four the evolving local bound of Theorem~1.1(3) is \(\Lambda_1(\epsilon+\sqrt{t})\). Curvature is bounded on each positive-time
subinterval and
\[
 \sup_{x,\ 0<t\le t_*}
 \|\Rm(t)\|_{L^2(B_t(x,1/2))}\le C\epsilon .
\]
Here and in choosing \(\epsilon\), the constants in this estimate
are independent of the chosen radius.

For clarity, the potential in the isotropic pair-sum evolution needs
control down to time zero. The weighted Sobolev estimate
\eqref{eq:smoothing-weighted-sob} follows from the initial Sobolev
bound and scalar lower bound exactly as in its proof. Those steps
use normalized test functions and hold at the rescaled volume
\(r^{-4}\), with both constants bounded independently of \(r\).
For a test supported on an evolving ball of radius \(1/2\),
H\"older bounds the scalar potential term by
\(C\epsilon\|f\|_4^2\). Absorption gives an ordinary local
Sobolev inequality with uniform coefficients. Iterating a ball
cutoff gives a uniform lower volume on every evolving ball of
radius \(1/16\). The total rescaled volume stays at most
\(e^{t_*}r^{-4}\), since the scalar curvature stays above \(-1\).

At each time take a maximal \(1/8\)-separated net. Disjoint small
balls bound its cardinality by \(N(r)\), uniformly in the index
and time. Use Lipschitz cutoffs on radius-\(1/4\) balls, equal
to one on the radius-\(1/8\) balls, and normalize by the square
root of their squared sum. The resulting functions satisfy
\(\sum_j\chi_j^2=1\) and \(\sum_j|d\chi_j|^2\le C(r)\).
Apply the local Sobolev estimate and the local curvature bound to
each \(\chi_j f\). The cross terms sum to zero, giving
\[
 \int |\Rm| f^2\,dV_t
 \le C\epsilon\int |df|^2\,dV_t
     +C(r)\int f^2\,dV_t .
\]
Only the zero-order coefficient depends on the covering number.
The cutoffs enter a spatial inequality at a fixed time.

For either negative isotropic pair sum \(w\) from
Section~\ref{sec:isotropic-gap}, test its weak evolution inequality
against \(2w\). Choose \(\epsilon\) once so the gradient term on
the right is absorbed. Since the scalar curvature is bounded
below, Gronwall gives
\(\|w(t)\|_2^2\le e^{C(r)t}\|w(0)\|_2^2=o_i(1)\) at this fixed
radius. The initial squared integral is scale invariant. At any
interior positive time the curvature and derivative bounds, injectivity
bound and total volume upper bound give a uniform diameter bound by packing along minimizing geodesics. Smooth compactness then gives a closed limit under
whole-manifold diffeomorphisms. Both pair sums are nonnegative there.
The positive-time strong maximum principle and \(b_2=0\) argument
in Corollary~\ref{cor:isotropic-ricci-tail} recognize its smooth
homotopy sphere as standard. This contradicts the nonstandard
types on the subsequence.

Now take a slow diagonal of fixed radii tending to zero. Transfer
the charged balls back to \(q_i\) by the distance and curvature
measure comparison, enlarging their radii by a factor tending to
one if necessary. Their volumes tend to zero by the assumed upper
bound. The modified scalar and spectral estimates imply vanishing
scalar and Weyl squared integrals there. The full tensor
decomposition then gives a definite trace-free Ricci integral.
Equation~\eqref{eq:modified-peak-energy} places their radii above
the maximum-curvature scale as claimed.
\end{proof}

Thus a sequence without shrinking balls carrying a fixed modified
Ricci charge cannot have uniformly bounded small-ball upper volume
ratios on any subsequence. The proposition uses neither a physical
chiral quantile nor an embedding of a compact exhaustion into a
charged ball.

\begin{theorem}[A nonflat flow at the modified charge scale]
\label{thm:charged-modified-flow}
Under the same upper volume bound, the radii can be chosen as the
first radii carrying a fixed small full curvature integral on the
smooth modified approximations. The corresponding pointed rescaled
Ricci flows converge on a subsequence, at positive times, to a
complete noncompact four-dimensional flow with nonnegative isotropic
curvature. On a ball of fixed radius about the marked center the
trace-free Ricci energy has a fixed positive lower bound at every
sufficiently small positive time. In particular the flow is nonflat.
\end{theorem}

\begin{proof}
Take the small norm threshold \(\epsilon\) in the preceding proof,
decreasing it below all the dimensional absorption thresholds below.
For each source the maximal ball integral is continuous in the radius
and attained in the center: distance spheres have zero volume, and
the source is smooth and compact. Its value at zero is zero. Let
\(a_i\) be the first radius at which its square root equals
\(\epsilon\), and choose an attaining center \(x_i\). The preceding
fixed-radius assertion gives \(a_i\to0\). Rescale by \(a_i^{-2}\)
and run the same closed flows, writing \(U=|\Rm|_{\rm tens}\).
Every initial unit ball now has curvature norm at most \(\epsilon\);
at \(x_i\) it is exactly \(\epsilon\).
The initial scalar curvature is at least \(-a_i^2\).

We need bounds uniform in this varying scale. Write
\(\|\Ric(t)\|_\infty\le b/t\) with \(b<1/2\), by the
curvature estimate in Lee--Lee Theorem~1.1(1), choosing \(\epsilon\)
and the common horizon sufficiently small. The shrinking balls
estimate gives \(B_t(z,1/2)\subset B_0(z,1)\) on that horizon
\cite[Lemma~3.2]{LeeLee2026Local}. Thus every evolving radius-\(1/2\)
ball has volume at most \(e^{t}V'\), from the original upper bound
and scalar lower bound, for a fixed \(V'\). Together with the
radius-\(1/16\) lower bound from the preceding proof, this bounds
the overlap of the radius-\(1/4\) cutoffs uniformly. For example,
the disjoint small balls about net points within \(1/4\) of a
point are all in its radius-\(1/2\) ball. The previous calculation
therefore gives uniform constants in the form bound
\[
 \int U f^2\,dV_t\le C\epsilon\!
       \left(\int |df|^2\,dV_t+\int f^2\,dV_t\right).
\]
The curvature tensor evolution, Kato's inequality and this bound
with \(f=U\) absorb the cubic term in the derivative of its
squared integral. The initial global full curvature budget then gives
\[
 \sup_{0\le t\le t_*}\int U^2\,dV_t+
 \int_0^{t_*}\!\int |\nabla\Rm|^2\,dV_t\,dt\le C_0 .
\]
The same form bound propagates vanishing of both negative pair sums
with a uniform Gronwall constant. With \(f=1\) it also gives
\(\int |R|\le C\operatorname{Vol}_t\); hence at every fixed time
the rescaled volumes tend to infinity. Curvature and injectivity
bounds, derivative estimates and pointed smooth compactness give
a complete limit flow on positive times. A compact smooth pointed
limit would identify its entire connected source along the
subsequence, contrary to the diverging volume. The propagated
pair-sum estimates give nonnegative isotropic curvature on the limit.

To locate the charge, fix a small terminal time \(\tau\).
All initial unit-ball points lie in \(B_\tau(x_i,D)\) for a
uniform \(D\). Indeed apply the local two-point distance estimate
of Lee--Lee Theorem~1.1 at a chain of centers along an initial
minimizing geodesic of length at most one. Its steps have length
less than the fixed comparison radius \(r_1/2\); each has uniformly
bounded length in the terminal distance. Centered versions of the
theorem concern the same smooth flow by compact uniqueness. Take
a terminal Lipschitz cutoff \(0\le\chi\le1\) equal to one on
this radius-\(D\) ball, zero outside radius \(D+1\), with
\(|d\chi|_{\tau}\le2\). The shrinking balls estimate puts its
entire support inside \(B_0(x_i,D+2)\). This fixed multiple of
the original \(a_i\) has vanishing unscaled modified volume.
Both original Weyl and scalar square integrals there tend to zero.

Normalize the Euler form as
\[
 \mathcal E=\left(|W^+|_{\rm HS}^2+|W^-|_{\rm HS}^2
                 +R^2/24-|Z|^2/2\right)dV .
\]
The decomposition \(U^2=4(|W^+|^2+|W^-|^2)+2|Z|^2+R^2/6\)
on the attaining ball now gives
\(\int\chi\mathcal E(0)\le-\epsilon^2/4+o_i(1)\).
In orthonormal bundles identified by the time-dependent isometries
of the Ricci flow, the time derivative of its metric connection
has norm at most \(C|\nabla\Ric|\): the differentiated Koszul
formula and the derivative of the moving frame each involve one
covariant derivative of \(\Ric\). Differentiating the Euler
polynomial gives an exact form with primitive bounded by
\(C|\nabla\Ric|U\) by the covariant Bianchi identity
\cite[\S3]{ChernSimons1974Characteristic}.
Keep the terminal cutoff fixed as a function of the source point.
Since \(|d\chi|_s\le2(\tau/s)^b\) for \(0<s\le\tau\), integration
by parts and the global curvature and derivative estimates give
\[
 \left|\int\chi\mathcal E(\tau)-\int\chi\mathcal E(0)\right|
 \le C C_0\left(\int_0^\tau(\tau/s)^{2b}\,ds\right)^{1/2}
 \le C'\sqrt\tau .
\]
One can first start at positive time and then use smooth continuity
at zero for each fixed source. Lipschitz tests are admissible in
the spatial integration by parts. The constants and support bound
are uniform in the source and in \(\tau\). Take \(C'\sqrt\tau\le
\epsilon^2/16\). The nonnegative Weyl and scalar terms in the terminal
Euler form force a definite trace-free Ricci integral on
\(B_\tau(x_i,D+1)\). Under complete pointed smooth convergence
every such ball is covered by a slightly larger compact comparison
at this fixed time. This transfers the lower bound to the limit on
a radius-\(D+2\) ball at every sufficiently small positive time.
The scalar curvature is nonnegative by the initial rescaled lower
bound. An interior zero would force it to vanish on all earlier
positive slices by the strong minimum principle. Its evolution
would give zero Ricci tensor there, contradicting the charge
just obtained. Thus the scalar curvature is positive on the flow.
\end{proof}

The model carries curvature after parabolic smoothing. The excluding
Euclidean-limit theorem uses first-order smooth markings for the
original rescalings at time zero. The distance comparison along the
flow alone does not furnish those markings.

\begin{corollary}[The initial end of the charged flow]
\label{cor:charged-flow-initial-end}
On the fixed manifold of the positive-time limit there is no sequence
\(t_j\downarrow0\) along which the metrics converge locally uniformly
and strongly in \(W^{1,4}_{\rm loc}\) to a smooth complete metric.
In particular the limiting flow has no smooth complete time-zero
extension. This assertion uses the markings at positive time only.
\end{corollary}

\begin{proof}
Pass to \(e_i\to e\), \(Y_i\to y>0\) and write
\(d=4e+y^2/12-8T<y^2/12\) for the limiting total modified
trace-free Ricci energy. Fix a common terminal time \(\tau_0\) small
enough for Theorem~\ref{thm:charged-modified-flow}. Write \(h(t)\)
for its limit flow on \(X\). Smooth pointed comparison uses the same
embeddings on closed positive-time subintervals in the flow gauge.
Transfer a fixed compactly supported smooth function \(f\) on \(X\)
to the sources via these embeddings, extending by zero. At \(\tau_0\)
its size, gradient and support radius are bounded uniformly in \(i\).
Its support lies in a fixed-radius initial ball in the rescaling by
the shrinking balls estimate. The rescaled volume there is uniformly
bounded at zero by the upper volume ratio and at all subsequent times
by the scalar lower bound. The rescaled initial scalar and Weyl
\(L^2\) norms on that ball tend to zero. For the transferred test
\(f_i\), \(|df_i|_s\le C_f (\tau_0/s)^b\).

Three estimates transport the initial data through smoothing. First,
\(\partial_t(R\,dV)=(\Delta R+2|\Ric|^2-R^2)dV\).
Testing against \(f_i\), the reaction terms have bounded spatial
integral, and integration by parts in the Laplacian costs at most
\(C_f(\tau_0/s)^b\|\nabla\Rm\|_2\). The initial scalar integral
against the test tends to zero by Cauchy--Schwarz. Sending \(i\) to
infinity at fixed time gives
\[
 \left|\int_X fR_{h(t)}\,dV_{h(t)}\right|
       \le C_f(t+t^{1/2-b}).
\]
Since \(h(t)\) has nonnegative isotropic curvature, both top Weyl
eigenvalues are bounded by \(R/6\). Trace freeness thus bounds both
Weyl norms by a multiple of \(R\). All three curvature terms vanish
locally in \(L^1\) as \(t\downarrow0\), using nonnegative cutoff
tests equal to one on the compact in question.

Second, the derivative of the Dirichlet energy of \(f_i\) is
the integral of \((2\Ric-Rg)(\nabla f_i,\nabla f_i)\).
Its absolute value is at most \(C_f s^{-2b}\) by the support-volume
and curvature bounds. The fourth-power volume integral changes
by at most \(C_ft\). At time zero the original rescaled Yamabe
inequality has constant \(Y_i\), and the scalar term against
\(f_i^2\) tends to zero. As \(2b<1\), both the quadratic energy
and fourth-power integral at time \(t\) differ from their initial
values by errors tending to zero as \(t\downarrow0\), uniformly
in \(i\). This transfers the full constant \(y\) to any proposed
locally uniform positive-definite time-zero limit.

Third, the Euler transport estimate with the gradient measured at
\(\tau_0\) bounds the change through time \(t\le\tau_0\) by
\(C_f t^{1/2-b}\). In the positive-time embeddings the initial
rescaled measures \(|Z_i|^2dV_i\) admit a local weak subsequential
limit \(\mu\), by compact exhaustion and the global bound.
It has total mass at most \(d\). In these same embeddings the
initial Euler forms tend weakly to \(-\mu/2\), since the initial
scalar and Weyl square integrals vanish on every compact support.
Consequently
\[
  \lim_{t\downarrow0}\int_X f\,\mathcal E(h(t))
       =-\tfrac12\int_X f\,d\mu .
\]
Moreover \(\mu\ne0\). The entire charged initial unit ball lies
in the fixed terminal ball \(B_{\tau_0}(x_i,D)\). A smooth limit
cutoff equal to one on a slightly larger closed terminal ball,
transferred to the sources, equals one there for large \(i\).
The initial trace-free Ricci integral there is at least
\(\epsilon^2/2-o_i(1)\). This argument fixes the terminal
marking before taking the time-zero trace.

Now suppose a sequence of times has the asserted limit \(h_0\).
The global curvature energy stays bounded. In fixed charts the
all-lowered curvature identity and strong first-order convergence
pass the curvatures to the smooth limit distributionally. The
local \(L^1\) bounds above give \(R_{h_0}=W_{h_0}=0\).
The transferred Yamabe test becomes
\(y\|f\|_{4,h_0}^2\le6\|df\|_{2,h_0}^2\).
Metric square roots identify the orthonormal bundles, with
connection differences tending to zero in \(L^4\). The Euler
transgression between each slice and \(h_0\) tends to zero on
compact smooth tests by the \(L^2\) curvature bound, as in
Theorem~\ref{thm:modified-shrinking-limits}. Thus the tested
trace gives \(|\Ric_{h_0}|^2dV_{h_0}=\mu\).

A complete such \(h_0\) has harmonic curvature, finite total
curvature and the homogeneous Sobolev inequality. Its finitely
many ALE ends and the end formula of
Theorem~\ref{thm:modified-shrinking-limits} apply. Indeed every
fixed smooth truncation of \(X\) already embeds in the original
homotopy spheres by a positive-time marking, which is all that
the topological part requires. The budget
\(\int|\Ric_{h_0}|^2\le d<y^2/12\) and the radial Yamabe test on
each end exclude two or more ends. The one-end Euler formula
then forces zero Ricci energy, contradicting \(\mu(X)>0\).
\end{proof}
